\section{A first example: the integers under addition}\label{sec:07-groups:03-integers-example:1}

A definition with no inhabitant is empty. Does anything satisfy \texttt{Group}? \((\mathbb{Z}, +, 0, -(-))\) does, and the construction proceeds field by field, each obligation closed by a short, explicit tactic proof rather than assumed.

\begin{lstlisting}
def intGroup : Group Int where
  op := fun a b => a + b
  id := 0
  inv := fun a => -a
  assoc := by
    intro a b c
    -- Goal: (a + b) + c = a + (b + c)
    exact Int.add_assoc a b c
  id_left := by
    intro a
    -- Goal: 0 + a = a
    exact Int.zero_add a
  id_right := by
    intro a
    -- Goal: a + 0 = a
    exact Int.add_zero a
  inv_left := by
    intro a
    -- Goal: (-a) + a = 0
    exact Int.add_left_neg a
  inv_right := by
    intro a
    -- Goal: a + (-a) = 0
    exact Int.add_right_neg a
\end{lstlisting}

Each field is a single \texttt{intro} followed by \texttt{exact}, naming the core-library lemma that already states the fact about \texttt{Int}. Nothing is hidden: \texttt{Int.add\_\allowbreak{}assoc}, \texttt{Int.zero\_\allowbreak{}add}, and the rest are themselves proved, elsewhere, by induction on the same representation of \texttt{Nat}/\texttt{Int} from Chapter 5, and reused here rather than re-derived.

\begin{mathreading}

This exhibits \((\mathbb{Z}, +, 0, -)\) as an object of \(\mathbf{Grp}\). \texttt{intGroup} is a proof that \(\mathbb{Z}\) is a group, built by supplying \((+, 0, -)\) and citing, not re-proving, each axiom: associativity is \(\mathrm{Int.add\_assoc}\), the identity laws are \(0 + a = a = a + 0\), the inverse laws are \((-a) + a = 0 = a + (-a)\). This is the one-line textbook remark ``\(\mathbb{Z}\) under addition is an abelian group,'' with the underlying inductions on \(\mathbb{Z}\) written out rather than assumed.

\end{mathreading}

\textbf{Mathlib equivalent.} Mathlib needs no \texttt{intGroup}-style bundle at all. \texttt{Int} is already an \href{https://loogle.lean-lang.org/?q=AddCommGroup}{\texttt{AddCommGroup}} instance, and the five axioms above are free-standing lemmas applying to every additive group, not only \texttt{Int}.

\begin{lstlisting}
example : AddCommGroup Int := inferInstance

example (a b c : Int) : (a + b) + c = a + (b + c) := add_assoc a b c
example (a : Int) : 0 + a = a := zero_add a
example (a : Int) : a + 0 = a := add_zero a
example (a : Int) : -a + a = 0 := neg_add_cancel a
example (a : Int) : a + -a = 0 := add_neg_cancel a
\end{lstlisting}

Same five facts about \(\mathbb{Z}\); where the book \emph{assembles} a \texttt{Group Int} term by hand, the instance already exists, found automatically by \href{https://loogle.lean-lang.org/?q=inferInstance}{\texttt{inferInstance}}. \href{https://loogle.lean-lang.org/?q=add_assoc}{\texttt{add\_\allowbreak{}assoc}}/\href{https://loogle.lean-lang.org/?q=zero_add}{\texttt{zero\_\allowbreak{}add}}/etc. are generic over \emph{any} \texttt{AddCommGroup}, so they apply equally to the \texttt{perm3Group}-style example of Section 4 once phrased in the \texttt{Group}/\texttt{AddCommGroup} classes of Mathlib, as that section does next.
